\documentclass[11pt]{article}

\usepackage[T1]{fontenc}
\usepackage{lmodern}
\usepackage[a4paper,margin=29mm]{geometry}
\usepackage{amsmath,amssymb,amsthm,mathtools}
\usepackage{enumitem}
\usepackage{xcolor}
\usepackage[colorlinks=true,linkcolor=blue!55!black]{hyperref}
\usepackage[nameinlink,capitalize,noabbrev]{cleveref}

\allowdisplaybreaks
\setlength{\parindent}{0pt}
\setlength{\parskip}{0.5em}

\newtheorem{theorem}{Theorem}[section]
\newtheorem{lemma}[theorem]{Lemma}
\newtheorem{proposition}[theorem]{Proposition}
\newtheorem{claim}[theorem]{Claim}
\theoremstyle{remark}
\newtheorem{remark}[theorem]{Remark}

\newcommand{\E}{\mathbb E}
\newcommand{\Pp}{\mathbb P}
\newcommand{\one}{\mathbf 1}
\newcommand{\opt}{\operatorname{OPT}}
\newcommand{\cost}{\operatorname{cost}}
\newcommand{\divg}{\operatorname{div}}

\title{Tail-batched active chords:\\an $O(n^{1/4}\sqrt{\log(n+1)})$ attack on the asymmetric clairvoyance gap}
\author{Proof addendum}
\date{5 August 2026}

\begin{document}
\maketitle

\begin{abstract}
This note closes the square-root discrepancy bottleneck in the
$O(n^{1/3})$ proof.  The backbone is sampled at rate $\alpha p_i$, so it has
only $O(1+\alpha\mu)$ vertices while its circulation costs
$O(\log(n+1)/\alpha)$ times the posterior optimum.  All outside clients with
the same chord tail are probed in chord-head order.  A causal, marginally
unbiased controller chooses one or four copies of every activated chord.  It
has the two simultaneous properties
\[
  \E[A_i k_i]=2p_i
  \quad\text{and}\quad
  \sup_t\E\exp\!\left(c\left|
       \sum_{i\le t}(A_i k_i-2p_i)\right|\right)=O(1).
\]
Thus the auxiliary flow can still be fixed before probing, while the error of
one tail is subexponential rather than the square root of its probability
mass.  Errors belonging to distinct tails are independent, so the global
cyclic-interval discrepancy is
$O(\sqrt{|R|\log(n+1)}+\log(n+1))$.  Optimizing $\alpha$ gives a
$\widetilde O(\mu^{1/3})$ low-probability policy.  Probability thresholding
then gives $O(n^{1/4}\sqrt{\log(n+1)})$.
\end{abstract}

\section{Statement}

We use the stochastic directed-metric TSP model and the notation of the
supplied $n^{1/3}$ proof note.  In particular, $B=\E[\opt(A)]$,
$\mu=\sum_i p_i$, probing is free and remote, and an active client must be
visited immediately.

Put
\[
  L:=\max\{1,\log(n+1)\}.
\]

\begin{theorem}\label{thm:main}
Every instance on $n$ clients has a causal randomized policy of expected cost
\[
  O\!\left(n^{1/4}\sqrt L\right)\opt_{\rm post}.
\]
The public and internal randomness can be fixed by averaging if a
deterministic policy is desired.
\end{theorem}

The new ingredient is the following low-probability statement.

\begin{proposition}[thinned-backbone policy]\label{prop:low}
Let $W$ be a subinstance, let
\[
 p_*:=\max_{i\in W}p_i\le\frac14,
 \qquad \mu:=\sum_{i\in W}p_i,
 \qquad B:=\opt_{\rm post}(W).
\]
For every $0<\alpha\le1$, there is a causal policy of expected cost at most
\begin{equation}\label{eq:low-alpha}
 O\!\left(
   \frac L\alpha+\sqrt{p_*\mu}
   +\sqrt{\alpha\mu L}+L
 \right)B.
\end{equation}
Consequently, choosing
$\alpha=\min\{1,(L/\max\{\mu,L\})^{1/3}\}$ gives
\begin{equation}\label{eq:low-optimized}
 O\!\left(L+L^{2/3}\mu^{1/3}+\sqrt{p_*\mu}\right)B.
\end{equation}
\end{proposition}

\section{An unbiased mean-reverting multiplicity controller}

The point of the next lemma is the conjunction of marginal unbiasedness and
exponential tightness.  A usual bang--bang controller has the latter but not
the former; independent randomized multiplicities have the former but retain
square-root fluctuation.

\begin{lemma}[one-dimensional controller]\label{lem:controller}
Let $q_1,\ldots,q_T\in[0,1/4]$ and let $A_t$ be independent Bernoulli
variables of means $q_t$.  There is a causal rule which, after observing
$A_t=1$, chooses $K_t\in\{1,4\}$, such that
\begin{align}
  \E[A_tK_t]&=2q_t &&(t\in[T]),\label{eq:marginal}\\
  X_t&:=\sum_{j\le t}(A_jK_j-2q_j),\label{eq:state}\\
  \sup_{t\le T}\E e^{c|X_t|}&\le C.\label{eq:exp-moment}
\end{align}
for universal constants $c,C>0$.  Decisions for disjoint input sequences can
be made with independent randomness.
\end{lemma}

\begin{proof}
We give the distributional controller explicitly.  Suppose the law of
$X_{t-1}$ has already been determined.  Attach an independent uniform rank
inside every atom and let $D_t$ be the bottom one third of this ranked law.
On $A_t=1$, use $K_t=4$ when $X_{t-1}\in D_t$ and $K_t=1$ otherwise.  At a
split atom the rank supplies the required randomization.  Hence
\[
  \Pp(X_{t-1}\in D_t)=\frac13,
  \qquad \E[K_t\mid A_t=1]=\frac13\,4+\frac23\,1=2.
\]
The current activation is independent of the old state, which proves
\eqref{eq:marginal}.  It also gives $\E X_t=0$ for every $t$.

For completeness, we record the quantile-drift calculation.  We use the
following elementary rearrangement inequality.  If $Z$ has mean zero and
$H(Z)$ equals $4$ on the bottom third of the ranked law of $Z$ and $1$ on the
top two thirds, then, for $0<c\le1/32$,
\begin{equation}\label{eq:quantile-covariance}
 \E\bigl[(H(Z)-2)\sinh(cZ)\bigr]
 \le-\frac16\bigl(\E\cosh(cZ)-1\bigr).
\end{equation}
To verify it, split at zero and at the one-third quantile.  On the bottom
third the coefficient is $2$, and on its complement it is $-1$.
Monotonicity shows that the average of $\sinh(cZ)$ on the former set is no
larger than on the latter.  In the two cases in which the quantile lies to
the left or to the right of zero, use respectively
\[
 \sinh u\ge\cosh u-1\quad(u\ge0),
 \qquad
 -\sinh u\ge\cosh u-1\quad(u\le0),
\]
and use $\E Z=0$ for the part straddling zero.  After multiplying by the
set masses $1/3$ and $2/3$, this gives
\eqref{eq:quantile-covariance}; the constant $1/6$ leaves slack for a split
atom.

Apply this with $Z=X_{t-1}$, $q=q_t$, and
$D:=A_tH(Z)-2q$.  Conditional on $Z$,
\[
 \E[D\mid Z]=q(H(Z)-2),
 \qquad |D|\le4,
 \qquad \E[D^2\mid Z]\le17q.
\]
Taylor's theorem and
$\cosh(c(Z+u))\le e^{c|u|}\cosh(cZ)$ give
\begin{align*}
 \E\cosh(cX_t)
 &\le \E\cosh(cZ)
   +cq\E[(H(Z)-2)\sinh(cZ)]
   +\frac{17}{2}c^2qe^{4c}\E\cosh(cZ)\\
 &\le\left(1-\frac{cq}{12}\right)\E\cosh(cZ)
   +\frac{cq}{6}
\end{align*}
after decreasing the universal $c$ if necessary.  Starting from $X_0=0$,
induction gives $\sup_t\E\cosh(cX_t)\le2$.  Since
$e^{c|x|}\le2\cosh(cx)$, this proves \eqref{eq:exp-moment} with $C=4$.
Every choice uses only the old state, its known law, fresh internal
randomness, and the current activation, so it is causal.
\end{proof}

Two consequences will be used repeatedly.  First, for every $s,t$,
Cauchy--Schwarz and \eqref{eq:exp-moment} give
\begin{equation}\label{eq:increment-exp}
  \E\exp\!\left(\frac c2|X_t-X_s|\right)\le C.
\end{equation}
Second, $\E X_t=0$.  Thus every interval increment of one controller is a
centered subexponential random variable with universal parameters.

\section{A thinned posterior backbone}

\begin{lemma}[thinned certificate]\label{lem:thin}
For every $0<\alpha\le1$ there are a fixed backbone
$R=S\cup\{r\}$, a rooted tour $T_R$, and a fractional circulation $z$ of the
same form as in the $n^{1/3}$ note such that, writing $m=|R|$,
\begin{align}
 \cost(T_R)&\le O(B),\label{eq:tour-cost}\\
 \cost(z)&\le O(LB/\alpha),\label{eq:z-cost}\\
 m&\le O(1+\alpha\mu).\label{eq:m-size}
\end{align}
Every outside client $i$ has incoming and outgoing $z$-degree $p_i$, and
there are no arcs between two outside clients.
\end{lemma}

\begin{proof}
Sample $S$ with independent probabilities $\alpha p_i$, independently of a
posterior sample $A$.  Put $C=A\cup S$, take a tie-broken optimal tour on
$C\cup\{r\}$, and use the padded-segment circulation construction from the
$n^{1/3}$ note.  Conditional on $C$, a vertex of $C$ is in $A\setminus S$
with probability
\[
 \frac{p_i(1-\alpha p_i)}{p_i+\alpha p_i-\alpha p_i^2}
 =\frac{1-\alpha p_i}{1+\alpha-\alpha p_i}
 \le\frac1{1+\alpha}.
\]
The classifications are independent.  Hence the longest outside run has
conditional expectation $O(L/\alpha)$.  Averaging the padded circulation
over $A$ therefore gives
\[
 \E_S\cost(z^S)
 \le O(L/\alpha)\E\opt(A\cup S)
 \le O(LB/\alpha),
\]
because $S$ is stochastically dominated by a posterior sample and
$\opt(A\cup S)\le\opt(A)+\opt(S)$.  Also
$\E\opt(S)\le B$ and $\E(|S|+1)=1+\alpha\mu$.  Apply averaging to the sum of
these three nonnegative quantities after their respective normalizations and
fix a suitable outcome $S$.
\end{proof}

Fix the certificate.  Let $U=W\setminus S$.  Couple flow through every
$i\in U$ into $x_{i,ab}$, $a,b\in R$, and let $y$ be the backbone-to-backbone
part, exactly as in the $n^{1/3}$ note.  Independently sample one public pair
$P_i=(a_i,b_i)$ with probability $x_{i,ab}/p_i$.  Put
\[
 b(P):=\sum_{i\in U}p_i(e_{a_i}-e_{b_i}),
 \qquad
 \bar b:=\E_Pb(P).
\]
Then $\divg(y)=-\bar b$, expected sampled-detour cost equals the detour part
of $z$, and the interval maximal inequality gives
\begin{equation}\label{eq:endpoint-error}
 \E_P\max_I|2(b(P)-\bar b)(I)|
 \le32\Bigl(\sum_{i\in U}p_i^2\Bigr)^{1/2}
 \le32\sqrt{p_*\mu}.
\end{equation}

Apply the integer interval-balancing lemma from the $n^{1/3}$ note to
$2y$, target vector $2b(P)$, and cycle $T_R$.  It gives a public nonnegative
integer flow $N(P)$ satisfying
\begin{align}
 \cost(N(P))
 &\le2\cost(y)+
 \left(2\max_I|(b(P)-\bar b)(I)|+1\right)\cost(T_R),
 \label{eq:N-cost}\\
 \max_I|(2b(P)+\divg N(P))(I)|&\le1.
 \label{eq:N-balance}
\end{align}

\section{Tail batching and global interval discrepancy}

For each $a\in R$, order all outside clients with $a_i=a$ by the cyclic
position of $b_i$; break ties publicly.  Run an independent copy of
\cref{lem:controller} on this list, with $q_t=p_i$.  When $i$ is probed and
active, the controller chooses $k_i\in\{1,4\}$.  If it is inactive, put
$k_i=0$.  For a head prefix $J$ define the row error
\[
 X_a(J):=\sum_{\substack{i:a_i=a\\b_i\in J}}
          (A_i k_i-2p_i),
\]
where a cut through a tied group uses the fixed tie order.  For a general
head interval use the corresponding difference of two prefixes.

\begin{lemma}[row aggregation]\label{lem:rows}
For the final selected chord multiset,
\begin{equation}\label{eq:row-bound}
 \E\max_{I\text{ cyclic interval}}
 \left|\sum_{i\in U}(A_i k_i-2p_i)
                  \chi_I(a_i,b_i)\right|
 \le O(\sqrt{mL}+L).
\end{equation}
Moreover, for each $i$,
\begin{equation}\label{eq:two-p}
 \E[A_i k_i]=2p_i.
\end{equation}
\end{lemma}

\begin{proof}
Equation \eqref{eq:two-p} is \eqref{eq:marginal}.  Fix a cyclic interval
$I\subseteq R$.  Put
\begin{equation}\label{eq:row-decomp}
 \Delta_I:=\sum_{a\in I}X_a(R)-\sum_{a\in R}X_a(I).
\end{equation}
The variables belonging to different tails $a$ are independent.  Every
summand in either sum in \eqref{eq:row-decomp} is centered, and
\eqref{eq:increment-exp} says that it has a universal subexponential norm.
Bernstein's inequality therefore gives, for every fixed $I$ and $u\ge0$,
\[
 \Pp\left(|\Delta_I|>u\right)
 \le4\exp\left[-c_0\min\left\{\frac{u^2}{m},u\right\}\right].
\]
There are at most $m^2$ cyclic intervals.  A union bound and tail integration
give \eqref{eq:row-bound}; replacing $\log(m+1)$ by $L$ only weakens the
bound.
\end{proof}

For the final public-and-realized chord multiset
\[
 F:=N(P)\cup\{k_i\text{ copies of }a_i\to b_i:A_i=1\},
\]
\eqref{eq:N-balance} and \cref{lem:rows} imply
\begin{equation}\label{eq:Lambda}
 \E[\Lambda(F)\mid P]
 \le O(\sqrt{mL}+L).
\end{equation}

\section{Causal execution and cost}

At a virtual backbone vertex $a$, use the following priority rule.
\begin{enumerate}[leftmargin=2em]
 \item Use an unused copy of a previously revealed selected chord leaving
 $a$, if one exists.
 \item Otherwise probe the next client in the head-ordered list at $a$.
 Inactive probes do not move.  If the client is active, run the controller,
 create $k_i$ copies, serve $i$ immediately along the first copy, and end the
 virtual detour at $b_i$.
 \item Once the list is exhausted, use an unused auxiliary arc of $N(P)$
 leaving $a$, if one exists.
 \item If none of the preceding actions applies, advance on the backbone
 cycle.
\end{enumerate}
Probe a backbone client at its first virtual visit.  Maintain the actual
position and pending-walk invariant from the causal-virtual-execution lemma in
the $n^{1/3}$ note.

This is a causal stack walk on the final multiset $F$.  Although a selected
chord is revealed only when its client is probed, a fallback cycle edge out of
$a$ is not used until the whole list at $a$ and every revealed outgoing copy
have been exhausted.  Therefore the proof of the active-first stack lemma
applies verbatim: if a cycle edge is the last one used for the first time,
then at that moment all lists and all chord stacks have been exhausted.
Consequently every backbone edge is used at most $\Lambda(F)+1$ times, plus
one final scan for residual inactive probes.

For every realization the actual cost is at most
\begin{equation}\label{eq:virtual-cost}
 \cost(N(P))+
 \sum_{i\in U}A_i k_i
       \bigl(d(a_i,i)+d(i,b_i)\bigr)
 +(\Lambda(F)+2)\cost(T_R).
\end{equation}
Take expectation.  By \eqref{eq:two-p}, the expected detour term is exactly
twice the sampled weighted detour cost.  Average \eqref{eq:N-cost} over $P$
and use \eqref{eq:endpoint-error}; then use
\eqref{eq:Lambda}, \eqref{eq:tour-cost}, \eqref{eq:z-cost}, and
\eqref{eq:m-size}.  This gives
\begin{align*}
 \E\cost(\text{policy})
 &\le O\left(\cost(z)+
   \bigl(\sqrt{p_*\mu}+\sqrt{mL}+L\bigr)\cost(T_R)\right)\\
 &\le O\left(\frac L\alpha+\sqrt{p_*\mu}
   +\sqrt{\alpha\mu L}+L\right)B.
\end{align*}
This proves \cref{prop:low}.  Notice that every multiplicity is selected at
the instant its activation is revealed; no posterior information is used.

\section{Threshold optimization}

We finish the proof of \cref{thm:main}.  Put
\[
 \tau:=\min\left\{\frac14,(nL^2)^{-1/4}\right\},
 \qquad
 H:=\{i:p_i\ge\tau\},
 \qquad
 W:=\{i:p_i<\tau\}.
\]
The fixed-tour binning lemma from the $n^{1/3}$ note gives a causal policy on
$H$ of expected cost at most
\[
 O(1/\tau)\opt_{\rm post}(H)
 =O(n^{1/4}\sqrt L)\opt_{\rm post}.
\]
For $W$, $p_*\le\tau$ and $\mu\le n\tau$.  Apply
\eqref{eq:low-optimized}:
\begin{align*}
 L^{2/3}\mu^{1/3}
 &\le L^{2/3}(n\tau)^{1/3}
 =O(n^{1/4}\sqrt L),\\
 \sqrt{p_*\mu}
 &\le\sqrt{n\tau^2}
 =O(n^{1/4}/\sqrt L).
\end{align*}
The additive $L$ is smaller after adjusting the universal constant.  Execute
the two policies successively from and back to the depot.  Monotonicity of
the posterior tour optimum and concatenation prove the theorem.

\section{What changed relative to the previous proof}

The old policy fixes one copy of every realized chord, so one tail containing
probability mass $M$ contributes a Bernoulli error of order $\sqrt M$.
Here all chords with that tail form a single causal controller.  The
controller reuses an active chord either once or four times; its marginal is
exactly two copies, which is why the deterministic auxiliary flow remains
available before probing.  The error of the whole tail is $O(1)$ at any fixed
head interval and $O(\log n)$ uniformly over its scan.  Only the number of
tails, namely the backbone size, now enters through a square root.  Thinning
the backbone is therefore useful: the two polynomial terms are
\[
  \frac L\alpha
  \qquad\text{and}\qquad
  \sqrt{\alpha\mu L},
\]
whose balance is $L^{2/3}\mu^{1/3}$.  This is precisely the leverage supplied
by repeated activated arcs.

\end{document}
